\begin{abstract}
How many characters can a language model simultaneously track in a narrative?
We present a systematic study that varies the number of characters from 2 to 64 in controlled synthetic stories, measuring how accurately frontier models report each character's current location and possession after a series of state changes.
We test two models---\gptfull and \gptmini---across three experimental conditions: a main scaling experiment, an extended range experiment (up to 64 characters), and a stress test with increased state-change density.
Both models achieve perfect accuracy up to 16~characters and maintain above 93\% accuracy even at 64~characters, far exceeding previously reported limits.
When errors do occur, they are overwhelmingly \emph{confusion errors}---the model assigns an attribute belonging to a different character---rather than forgetting errors.
For \gptmini, 93.9\% of all errors are confusions.
This error pattern suggests that character attributes are stored in overlapping distributed representations rather than independent memory slots, consistent with recent mechanistic interpretability findings on transformer state tracking.
We also find, counterintuitively, that \gptmini outperforms \gptfull, and that degradation is gradual rather than cliff-like.
Our results establish a new empirical baseline for entity tracking capacity in frontier models and provide error-pattern evidence bearing on the internal mechanisms transformers use for multi-entity reasoning.
\end{abstract}
